
\section{Program input}
\label{sec:program-loaders}
The programs in this section are designed to make it easier to load programs into the LGP-21 than tediously typing them on the Flexowriter and depositing them word-by-word into memory.
\subsection{Program Loader 1}
\subsubsection*{J1-10.0}
\label{sec:PIRloader}
\subsubsection*{PURPOSE}

\emph{Program Input 1}, J1-10.0, [Translator's note: this program is known more colloquially as ``PIR''.] is a memory program. It facilitates the reading of hexadecimal tapes—whether arbitrarily relocated or not—and calculates a checksum during the reading process. It can also read and check arbitrarily relocated decimal programs.

Program J1-10.0 performs the following functions:
\begin{enumerate}
\item Reading, encoding, and storing instruction words to predefined memory locations.
\item Reading and storing hexadecimal or alphanumeric words to predefined memory locations without encoding.
\item Setting an address modifier in a program to a specific value.
\item Reading decimal addresses incremented by the modifier value.
\item Setting the ``Start Fill'' counter to a specific address value.
\item Reading specific instructions and executing them immediately. 
\item Reading decimal addresses, followed by a jump either to this memory location or to the memory location specified by the address plus modifier.
\item Reading a decimal address; the contents of this memory location appear in the accumulator.
\item Reading, encoding, and storing arbitrarily relocatable decimal tapes.
\item Reading and storing hexadecimal tapes.
\item Selecting a specific input unit.
\end{enumerate}
Hexadecimal tapes are usually generated by a hexadecimal dump program, such as the J4-10.0 program.

\subsubsection*{Memory usage}
\begin{center}
\begin{tabular}{ll}
    \toprule
    Hexadecimal input only & 2 tracks \\
    Combined input & 3 tracks, 29 sectors \\
    Buffer usage & None \\
    \bottomrule
\end{tabular}
\end{center}textsf

\begin{itemize}
\item \texttt{MNNKHHHH}
\label{sec:hexadecimal-fill}
This keyword reads \texttt{NN} hexadecimal words in ascending order into consecutive memory locations, starting at memory location \texttt{HHHH} plus the modifier. The \texttt{K} is a store instruction. The input keyword contains a flag bit (bit 31); 
each of the NN words that are to be modified must also contain a flag bit.

%% CLARIFY: where us the sign bit? How is this specified? is the flag bit" the sign bit?

The address modifier (see below) is added to the operand address portion of each flagged word before it is stored.

A checksum is calculated from all \texttt{NN} words as they are being stored plus the input keyword (all values are checksummed without the address modifier). If the calculated checksum does not match the checksum on the tape following these words, an error stop occurs. If the checksums match, the computer requests a new input keyword.

Example: Let the modifier be $1200$ (decimal) and the input keyword be M40K1201' (hexadecimal). The next 64 (40 hex) hexadecimal words are read and modified (if they have a flag bit) and are stored consecutively, beginning in memory location $3000$.

\item \texttt{VNNCHHHH} - 

The hexadecimal fill keyword for non-relocatable [absolute] tapes reads in \texttt{NN} hexadecimal words starting with memory location \texttt{HHHH}. \texttt{C} is a store instruction. The input keyword does not contain a flag bit. The most recently-entered address modifier value is added to the operand address portion of all hexadecimal fill words that have a flag bit as they are stored.
%% Again, clarify where the fuck this flag bit is coming from on the command!

A checksum is calculated from all \texttt{NN} words as they are being stored plus the input keyword (all values are checksummed without the address modifier). If the calculated checksum does not match the checksum on the tape following these words, an error stop occurs. If the checksums match, the computer requests a new input keyword.

%% VERIFY: no example for this one?

\item \texttt{VNNCTTSS} 

Example: The input keyword \texttt{V40C1100'} causes the following 64 (40 hex) hexadecimal words to be stored, beginning in memory location $1700$.

When reading hexadecimal tapes with \texttt{V} or \texttt{M} keywords, the Program Input Routine verifies each stored word; specifically, after a word is stored, the value in memory is compared with the read-in value. If the two values do not match, an error stop occurs.

%% verify no example here too

\item \texttt{/000TTSS}

The Modifier Input Keyword sets the address modifier to the address \texttt{TTSS} contained within the input keyword.

When entering hexadecimal tapes, this Address Modifier is added to all keywords and words that contain a flag bit.

When entering decimal instruction words, the modifier is added to the operand addresses of all instructions except those preceded by an ``X''.

The modifier remains unchanged until it is replaced by a new Modifier Input Keyword.

Example: The input keyword \texttt{/0001200'} causes $1200$ to be added to all flagged operand addresses and keywords.

\item \texttt{.000TTSS'}

  The Absolute Stop and Jump Input Keyword causes the computer to stop. If the Start button is subsequently pressed, a jump occurs to cell \texttt{TTSS}. If the \textsf{PS-32} switch is set to \textsf{ON}, execution continues at the target address.

    Example: The input keyword $.0001663'$ would cause the computer to jump to $1663$ and halt execution (if \textsf{PS-32} is off).

\item \texttt{.100TTSS'}

  The Relative Stop and Jump Input Keyword causes the computer to stop. If the Start button is subsequently pressed, a jump occurs to cell \texttt{TTSS} plus the modifier. If the \textsf{PS-32} switch is set to \textsf{ON}, the stop is suppressed.

Example: If the modifier has been set to $1200$, then the input keyword \texttt{.1001663'} would cause a jump to $2863$.

\item \texttt{S000TTyy'}

The Selection Keyword selects the specific input unit designated by the decimal track address \texttt{TT}. (\texttt{yy} can be any arbitrary sector address; it is ignored, but must be supplied). A list specifying the corresponding input units for specific track addresses can be found in the LGP-21 Programming Manual.

%% ADD reference for that manual.

After entering the Selection Keyword, the computer stops. After pressing the Start button, the computer continues reading via the selected input unit. 

Restarting J1-10.0 by jumping to the start of the program will automatically reselect the primary Flexowriter for input. Reading will then continue via the Flexowriter until a new Selection Keyword selects a different unit.

Example: If 32 refers to a second Flexowriter, then the keyword \texttt{S0003200'} selects this second typewriter for further input.

\end{itemize}
\subsubsection*{Decimal input}

\begin{itemize}
\item \texttt{000TTSS'}

The Decimal Fill keyword causes the input program to store decimal words in ascending order into consecutive cells, beginning with memory location \texttt{TTSS}.

Any valid address followed by a stop code may be used. Every decimal input should be preceded by a Decimal Fill keyword. The ``start fill'' address is stored in the start-fill counter. This memory location of the subroutine contains the address for the decimal input. The counter is incremented by 1 each time after a word (not a keyword) is read and stored. The read-in words are stored consecutively until a new decimal fill keyword is read or the input is terminated.

Example: The fill keyword \texttt{;0000400'} would store read-in words beginning with memory location 0400.

\item \texttt{,00000NN'}

The keyword for entering hexadecimal words causes \texttt{NN} hexadecimal words to be stored, specifically beginning in the memory cell whose address is held in the start-fill counter. \texttt{NN} must be specified in decimal and can take a value from 1 to 63. The stored words are not incremented by the address modifier. Leading zeros may be omitted for the hexadecimal words that are to be read in.

Each hexadecimal word and keyword must be followed by a stop code. If a zero is to be entered, it is not necessary to write a \texttt{0}. A stop code alone is sufficient to indicate a zero. Each hexadecimal word must not consist of more than 8 characters.

Example: The input keyword \texttt{,0000003'} followed by the three hexadecimal words \texttt{-2'JK73''} would cause the values \texttt{00000002}, \texttt{0000JK73}, and \texttt{00000000} to be stored into the next three consecutive cells.

\item \texttt{A00000NN'}

The keyword for entering alphanumeric words causes \texttt{NN} alphanumeric words to be stored, beginning in the memory cell whose address is held in the start-fill counter. The alphanumeric words are shifted two places to the left and stored without binary conversion. \texttt{NN} must be specified in decimal and can take a value from 1 to 63.

Each word can contain up to five characters in 6-bit representation (if more than five characters per word are entered, only the last five are retained). The last stored word receives a group-end indicator, i.e., $1@30$. All characters except tape feed (blank tape), delete code, and stop code can be entered; for example, lowercase letters can be stored.

Example: The input keyword \texttt{A0000004'} followed by the four words \texttt{TOTAL' UNIT'S PRO'DUCED'} would cause the information \texttt{TOTAL UNITS PRODUCED} to be stored into four consecutive memory locations without conversion into binary.

\item \texttt{+00LTTSS'}

The Instruction Input Keyword is treated by the subroutine exactly like an instruction.

The arbitrary instruction \texttt{LTTSS} contained within the keyword—where \texttt{L} is any positive machine instruction and \texttt{TTSS} represents any valid address—is executed immediately after either the input of a subsequent hexadecimal word or the input of a stop code (if no subsequent word is required). The address portion of the input keyword is not altered by the Address Modifier.

Each hexadecimal keyword and input word must be followed by a stop code. Left-side leading zeros within a word may be omitted. The computer returns to an input command without the execution of the instruction causing a branch. Jump instructions should be avoided.

Note: The instruction contained within the input keyword does not alter the fill counter.

Example: The input keyword \texttt{+00h1637'} followed by the hexadecimal word \texttt{73W08'} would cause the value \texttt{00073W08} to be stored in memory cell 1637.

\item \texttt{-000TTSS'}

The Sequential Display Input Keyword causes the contents of memory cell \texttt{TTSS} to appear in the accumulator when the computer stops. A start signal causes a return to an input command for a new keyword, provided that the \texttt{PST} switch is off. If the \texttt{PST} switch is on, the contents of \texttt{TTSS$ + i$} (where $i = 1, 2, \dots$) will appear in the accumulator.

With each start signal, $i$ is incremented by 1. It is also possible to obtain the address of the displayed word by switching to single-step operation and starting the machine. (The \texttt{B} instruction that fetched the information just displayed will appear in the accumulator.) \texttt{TTSS} may be any arbitrary valid address.

Example: The input keyword \texttt{-0003263'} would directly display the contents of memory cell $3263$ in the accumulator.

Every entered word whose leftmost four bits correspond either to an ``8'' (negative instruction) or a zero (positive instruction) is treated as an instruction. The operand address is altered by the Address Modifier, unless the instruction is preceded by an ``X''. The subroutine allows the inclusion of index tags (used by some interpretive systems) immediately preceding the instructions. The inclusion of an ``X'' or an index tag does not alter the instruction in memory. If the leftmost four bits of a word correspond to any of the numbers $9$ through $13$, an error stop occurs (see the list of error stops).

Example: Let the Address Modifier be $0600$. Different types of instruction words would be stored as follows:

\begin{tabular}{lll}
Input     &   Storage \\
B 1234    &    B 1834 \\
XB 1234   &    B 1234 \\
800T 1234 &    800T 1834 \\
80XT 1234 &    800T 1234 \\
8B 1234   &    8B 1834 \\
X8B 1234  &    8B 1234 \\
\end{tabular}

The subroutine consists of two parts: hexadecimal input and decimal input. The hexadecimal input section can be used independently of the decimal input; however, the decimal input is dependent on the hexadecimal section, which must always be stored in memory.
\end{itemize}

\subsubsection*{CALLING SEQUENCE}
The program is called manually by jumping to the starting address (regardless of which section is being used).

\subsubsection*{OPERATING INSTRUCTIONS}

PROGRAM INPUT
The tape contains the program in relocatable hexadecimal form with its own bootstrap loader. Operation is performed as follows:

\begin{enumerate}[label=\alph*)]
\item Place the program tape into the typewriter's reader mechanism and release all levers.
\item \label{step:input} Set the mode selector switch to \texttt{Manual Input} on the console.
\item Press \texttt{Start Read} on the Flexowriter.
\item Press \texttt{Fill Clear} on the console.
\item Press \texttt{Start Read}.
\item If the program is to be relocated, enter the starting address in hexadecimal form via the typewriter (typically the starting address is $0000$); otherwise, omit this step.
\item Set the mode selector switch to \texttt{Single Step}.
\item \label{step:run} Press \texttt{Execute} on the console.
\item Repeat steps \ref{step:input} through \ref{step:run} until ``normal'' is completely printed out by the computer.
\item Set the mode selector switch to \texttt{NORMAL}.
\item Press \texttt{START} on the console.
\item The program automatically continues reading in.
\end{enumerate}
Note: If the \texttt{PS-32} switch is turned \emph{off}, the computer stops after reading in the hexadecimal input section. If the \texttt{PS-32} switch is turned \emph{on} or if \texttt{START} is pressed, the decimal input section or a hexadecimal object program tape will be read in.

\texttt{PS-32} should be turned \emph{off} before the decimal input section is completely read in; otherwise, the reader mechanism will continue to operate without tape after the read-in process is complete.

The bootstrap always occupies memory cells $0000$ and $0002\dots0008$ (absolute), while the remainder resides in the area designated for the hexadecimal input section.

\subsubsection*{INSTRUCTIONS FOR RELOADING}

The bootstrap remains intact in memory if the program has not been relocated. To re-load the program without the bootstrap, proceed as follows:
\begin{enumerate}[label=\alph*)]
\item Place the program tape into the reader, specifically aligning it at the blank tape section between the bootstrap and the program.
\item Set the mode selector switch to \texttt{Manual Input} on the console.
\item Enter \texttt{C0000} via the typewriter.
\item Press \texttt{Fill Clear} on the console.
\item Enter the modifier in hexadecimal form (all 8 characters), or enter 8 zeros if the program is not to be relocated.
\item Set the mode selector switch to \texttt{Single Step} on the console.
\item Press \texttt{Execute} on the console.
\item Set the mode selector switch to \texttt{Manual Input} on the console.
\item Enter \texttt{U0008}.
\item Press \texttt{Fill Clear} on the console.
\item Set the mode selector switch to \texttt{Single Step} on the console.
\item Press \texttt{Execute}.
\item Set the mode selector switch to \texttt{NORMAL} on the console.
\item Press \texttt{Start}. The program will be read in automatically.
\end{enumerate}

\subsubsection*{The \texttt{C0000} and \texttt{U0008} Instructions}
In the LGP-21 instruction set, \texttt{C} is the Clear and Transfer instruction (storing the accumulator contents in memory), and \texttt{U} is the Unconditional Jump instruction. Entering \texttt{U0008} manually forces the to skip the initial bootstrap block and jump straight to the start of the primary loader routine at location 0008.

\subsubsection*{REQUIRED INPUT/OUTPUT UNITS}
Input units other than the standard typewriter must be specified by a selection code. Output units are not required.

\subsubsection{THE EFFECT OF THE PROGRAM JUMP SWITCHES}
\begin{itemize}
\item \texttt{PS-32} (Program Switch 32)
\begin{itemize}
\item \texttt{OFF}: 
\begin{itemize}
    \item The computer stops before executing a jump when a Stop and Jump Input Keyword is used.
    \item The computer stops after reading in the hexadecimal input section of the subroutine.
\end{itemize}
\item \texttt{ON}
\begin{itemize}
    \item No stop occurs when a Stop and Jump Input Keyword is used.
    \item No stop occurs after reading in the hexadecimal or decimal section of the subroutine.
\end{itemize}
\end{itemize}
\item \texttt{PST} (Program Switch-Typewriter)
\begin{itemize}
\item \texttt{OFF}: Only the contents of a single, specific memory cell will appear in the accumulator.
\item \texttt{ON}: Consecutive memory cells can be loaded into the accumulator by repeatedly pressing the Start button.
\end{itemize}
\end{itemize}

\subsubsection*{MESSAGES}
\subsubsection*{E (Input Error)}
Cause:
\begin{itemize}
    \item Hexadecimal Tape Input: This character indicates a checksum error.
    \item Decimal Tape Input: This character indicates an invalid input code.
\end{itemize}
Remedy:
\begin{enumerate}[label=\alph*)]
\item Move the tape back to the last M or V input keyword.
\item Press START (the input unit does not need to be re-selected).
\end{enumerate}
Remedy for an Invalid Input Code (the last word read contains the error):
\begin{enumerate}[label=\alph*)]
\item Press Manual Input on the typewriter.
\item Set the mode selector switch to Manual Input.
\item \label{step:enter-jump} Enter a jump to the starting address of the subroutine in hexadecimal form (i.e., U0000 or the respective starting address). If the subroutine begins in memory location 0000, steps \ref{step:enter-jump}, \ref{step:single-op}, and \ref{step:execute} may be omitted. 
\item Press Clear and Input (Fill Clear) (TR).
\item \label{step:single-op} Set the mode selector switch to Single Operation (Single Step) (TR).
\item \label{step:execute} Press Execute (TR).
\item Set the mode selector switch to Normal (TR).
\item Press Start (TR). The typewriter lamp lights up.
\item Input the correct input code.
\item If the typewriter is the active input unit, release the Manual Input lever, and the tape will continue reading.
\item If another input unit is used, enter its selection code, press the Computer Start lever, then press Start (TR), and the tape will continue reading.
\end{enumerate}

\subsubsection*{W (Write/Verify Error)}
This printout indicates that a verification error occurred during the reading of a hexadecimal tape (i.e., the word actually stored in memory did not match the word read from the tape). 

This error is remedied in the same manner as a checksum error.

\subsubsection*{Tape verification}
By modifying two instructions in the subroutine, it can be used to verify hexadecimal tapes. For example, after a hexadecimal tape has been punched, alter the two instructions in the subroutine and then read the hexadecimal tape in.

The subroutine will read the tape (without storing it) and compare it word-for-word with memory. The checksums are also verified. If the comparison is unsuccessful, the program halts with the error stop ``W'' or ``E''. The tape is correct if it reads through without difficulty.

The instructions to be modified are:

\begin{tabular}{llll}
 & Memory Location  &  Old Contents &   New Contents \\
 & L + 24     &        H J      &       U 0025 \\
 & L + 31     &        Y 0024   &       U 0032
\end{tabular}

% =========================================================================
% SECTION: QUICK REFERENCE & NOTATION LEGEND
% =========================================================================
\subsubsection*{Address Notation Legend}

\noindent
\hrule height \heavyrulewidth
\nopagebreak
\begin{description}[leftmargin=2.0em, itemsep=3pt, topsep=4pt, parsep=0pt, partopsep=0pt]
    \item[\texttt{Ø}] Machine instruction code, represented as a letter.
    \item[\texttt{TT}] Track portion of the address, in decimal (00--63).
    \item[\texttt{SS}] Sector portion of the address, in decimal (00--63).
    \item[\texttt{NN}] Number of words to read, in decimal.
    \item[\texttt{HHHH}] Hexadecimal memory address.
\end{description}
\nopagebreak
\hrule height \heavyrulewidth

\vspace{1.5em}

% =========================================================================
% SECTION: KEYWORD CONFIGURATIONS
% =========================================================================
\subsubsection*{Instruction Keyword Quick Reference}

\noindent
\begin{tabular}{ll}
\toprule
\textbf{Keyword Format} & \textbf{Basic Meaning} \\
\midrule
\texttt{ØTTSS}    & Instruction -- may be modified. \\
\texttt{XØTTSS}   & Instruction -- may not be modified. \\
\texttt{800ØTTSS} & Negative instruction -- may be modified. \\
\texttt{80XØTTSS} & Negative instruction -- may not be modified. \\
\texttt{IØTTSS}    & Instruction with index tag -- may be modified. \\
\texttt{XIØTTSS}   & Instruction with index tag -- may not be modified. \\
\bottomrule
\end{tabular}

% =========================================================================
% SECTION: DETAILED OPERATIONS SUMMARY (FOOLPROOF LONGTABLE)
% =========================================================================
\subsubsection*{Detailed Operation Summary}
\vspace{-1.2em} % <--- Add this line here to yank the table up tight against the header
\begin{longtable}{l p{11.5cm}}
\toprule
\textbf{Keyword} & \textbf{Operation \& Detailed Description} \\
\midrule
\endhead
\bottomrule
\endfoot

MNNKHHHH  & \textit{Fill with Relocatable Words} \newline
            Read NN hexadecimal words into memory, starting at location HHHH + modifier. If bit 31 is 1, operand address is modified. Checksum calculated and compared with word NN+1. \\
\addlinespace
VNNCHHHH  & \textit{Fill with Non-relocatable Words} \newline
            Read NN hexadecimal words into memory, starting at location HHHH. If bit 31 is 1, operand address is modified. Checksum calculated and compared with word following the NN words previously read. \\
\addlinespace
/000TTSS  & \textit{Set Modifier} \newline
            Sets an address modifier to TTSS. Modifier is added to all hexadecimal fields that have bit 31 set, and to the addresses of decimal instructions not preceded by an X. \\
\addlinespace
.000TTSS  & \textit{Stop and Jump Absolute} \newline
            Sets program counter to TTSS and stops execution. If PS-32 is ON, execution continues. \\
\addlinespace
.100TTSS  & \textit{Stop and Jump Relative} \newline
            Sets program counter to TTSS + modifier and stops. If PS-32 is ON, execution continues. \\
\addlinespace
S000TTxx  & \textit{Device Select} \newline
            Selects (decimal) input device TT (xx is ignored) and stops. To continue, press START. \\
\addlinespace
;000TTSS  & Decimal Fill \newline
            Reads decimal instructions into memory starting at address TTSS. \\
\addlinespace
,00000NN  & Hexadecimal Fill \newline
            Reads NN hexadecimal words into the following NN consecutive memory cells. \\
\addlinespace
A00000NN  & Alphanumeric Fill \newline
            Reads NN alphanumeric words into the following NN consecutive memory cells. \\
\addlinespace
+00ØTTSS  & Execute Immediate \newline
            Executes the instruction ØTTSS contained within the keyword. If the instruction has no operand, enter a stop code. Otherwise enter the operand word in hexadecimal notation an a stop code. \\
\addlinespace
-000TTSS  & Sequential Dump \newline
            The contents of memory cell TTSS are shown in the console's accumulator display when the computer stops. After pressing START, the computer requests an command if PST is OFF. If PST is ON, pressing START causes the contents of TTSS$ + 1$ to appear in the accumulator, and so on. The address of the cell currently shown in the accumulator display appears in the accumulator if you switch to ``Single Operation'' and press Start. \\
\end{longtable}

\newpage